\documentclass{article}
\usepackage[utf8]{inputenc}
\usepackage{xcolor}
\usepackage{hyperref}

\usepackage{rotating}

\usepackage{tabls}
\usepackage{booktabs}
\usepackage{amsmath}
\usepackage{amssymb}
\usepackage{amsthm}
\usepackage{amsfonts}
\usepackage{multicol}
\usepackage{enumitem}
\usepackage{microtype}
\usepackage{tikz}
\usepackage{pgfplots}
\usetikzlibrary{positioning}
\usepackage{multirow}

\usepackage{comment}

\usepackage{graphicx}
\usepackage{wrapfig}

\theoremstyle{plain}
\newtheorem{theorem}{Theorem}
\newtheorem*{theorem*}{Theorem}
\newtheorem{lemma}[theorem]{Lemma}
\newtheorem*{lemma*}{Lemma}
\newtheorem{proposition}[theorem]{Proposition}
\newtheorem{corollary}[theorem]{Corollary}

\theoremstyle{box}
\newtheorem{definition}[theorem]{Definition}
\newtheorem*{axiom*}{Axiom}
\newtheorem{dfn}[theorem]{Definition}
\newtheorem*{definition*}{Definition}
\newtheorem*{exercise*}{Exercise}
\newtheorem{observation}[theorem]{Observation}
\newtheorem{remark}[theorem]{Remark}
\newtheorem{example}[theorem]{Example}
\newtheorem{question}[theorem]{Question}
\newtheorem*{prep-problems}{Preparation Problems}

\newcommand{\pd}[3]{\frac{\partial^{#3} {#1}}{\partial {#2}^{#3}}}
\newcommand{\fpd}[2]{\frac{\partial {#1}}{\partial {#2}}}
\newcommand{\diff}[2]{\frac{d {#1}}{d {#2}}}

\begin{document}
\begin{definition*}[Describing Functions]\hfill
\begin{itemize}
\item We say a function $f$ is \underline{constant on the interval $I$} if for all $x_1, x_2$ in the interval, $f(x_1)=f(x_2)$.
\item We say a function $f$ is \underline{increasing on the interval $I$} if for all $x_1, x_2$ in the interval, $f(x_1) \leq f(x_2)$ when $x_1 < x_2$.
\begin{itemize}
\item We say a function $f$ is \underline{\textbf{strictly} increasing on the interval $I$} if for all $x_1, x_2$ in the interval, $f(x_1) < f(x_2)$ when $x_1 < x_2$.
\end{itemize}
\item We say a function $f$ is \underline{decreasing on the interval $I$} if for all $x_1, x_2$ in the interval, $f(x_1) \geq f(x_2)$ when $x_1 < x_2$.
\begin{itemize}
\item We say a function $f$ is \underline{\textbf{strictly} decreasing on the interval $I$} if for all $x_1, x_2$ in the interval, $f(x_1) > f(x_2)$ when $x_1 < x_2$.
\end{itemize}
\end{itemize}
\end{definition*}
%Definitions adapted from OpenStax Calculus Volume 1.

\vfill

\textbf{Exercises}: Practice with these definitions.
\begin{enumerate}
\item Draw a picture of a constant function.
\item Draw a picture of an increasing function.
\item Draw a picture of an increasing function that is not strictly increasing.
\item Draw a picture of a decreasing function.
\item Draw a picture of a decreasing function that is not strictly decreasing.
\item Write the equation of a constant function.
\item Write the equation of an increasing function.
\item Write the equation of a decreasing function.
\item In 3-5 sentences, describe a constant, decreasing, or increasing function (pick one). Give a specific example, tell a story.
%\item In 3-5 sentences, describe an increasing function (give a specific example, tell a story).
%\item In 3-5 sentences, describe a constant function (give a specific example, tell a story).
\end{enumerate}

\vfill

\textbf{Exercises}: Practice with R and these definitions.\\ 
Identify the intervals on which the function is increasing, constant, or decreasing.
(Adapted from Exercises 63-66 page 70, ``Mathematical Modeling and Applied Calculus" by Joel Kilty and Alex M. McAllister)
\begin{enumerate}
\item $f(x) = \begin{cases} 
      -2x + 1 & x < 5 \\
      4e^x & 5 \leq x
   \end{cases}.$
%\item $g(x) = \begin{cases} 
%      -2e^x & x < 0 \\
%      2x+7 & 0 \leq x
%   \end{cases}.$
\item $g(x) = \begin{cases} 
      10^{-2x} & x < -1 \\
      x + 4 & -1 \leq x
   \end{cases}.$
\item $h(x) = \begin{cases} 
      -3x -1 & x \leq 0 \\
      \frac{1}{2}(4^{-3x}) & 0 < x
   \end{cases}.$
\end{enumerate}

\end{document}